\begin{itemize}
\item \textbf{Scale and statistics.} Tiny models, \rhval{const/train_steps} steps, one training set size, \rhval{const/n_seeds} seeds and \rhval{const/n_probs} problems per task and seed. Paired $t$-tests with five seeds are underpowered for small effects and are uncorrected for multiple comparisons; ``not significant'' means ``not detected''.
\item \textbf{Synthetic, easy calculus.} The rules are invertible, so a wrong tactic never makes a goal unprovable; the policy only affects proof size. This is a much weaker notion of ``tactic prediction'' than in interactive provers with dead ends, and the conclusions may not transfer to non-invertible systems (intuitionistic logic, first-order logic with instantiation) or to real libraries. The task contains no arithmetic.
\item \textbf{Problem distribution.} Four hand-designed generators, rejection to a size range, 5 variables. Sizes of the test sequents are strictly larger than training, but the generators (e.g.\ substitution instances and rewrites) also shift structure with size, so ``size'' and ``structure'' are confounded. Optimal proof size was computed only for training sequents; we did not measure how far from optimal the test proofs of each system are.
\item \textbf{Baseline strength.} The heuristic was tuned over three variants on a bin from the target size regime; it is a hand-designed rule that exploits the known structure of G3cp, which is the strongest hand baseline we could write in the time available but is not provably optimal (the MLP with rank cost beats it). The policies were not tuned at all; a tuned policy (longer training, more capacity, other losses, value or proof-size heads, search with learned value) could behave differently.
\item \textbf{Post-hoc analysis.} The cost-function and MLP-rank results were run after the main runs revealed the weakness of $-\log p$ search; they are exploratory. We did not isolate \emph{why} $-\log p$ cost hurts.
\item \textbf{Reimplementation.} BFS, the heuristic search and both policies are our reimplementations in a shared engine; absolute numbers cannot be compared with published provers. Axiom closure is free and the first open goal is always expanded; other goal-selection and expansion-counting conventions would change absolute numbers.
\end{itemize}
